\section[PHT3D Exercise 5 Transport of uranium in a fractured rock aquifer]{PHT3D Exercise 5 Transport of uranium in a fractured rock aquifer}

In this exercise we will investigate the transport of solutes 
in a fractured aquifer system. The exercise will illustrate 
how the dual-domain mass transfer model (DDMT) can be used to mimic the
characteristic physical transport behaviour of fractured rock
and how the model mass transfer with the matrix may influence 
the transport of uranium. The illustrative example
is defined for a case in which the potential impact of 
the release of mining wastes on a surface water body 
is investigated. The main focus is on understanding the 
coupled physical and chemical transport behavior of uranium.        

\subsection{Groundwater flow}

The model for the site consists of a cross section that follows 
a principle flow path between the groundwater divide and 
a discharge point to a creek that is located 2 km downstream of the 
groundwater divide. Groundwater recharge from rainfall is the only 
driver of groundwater flow in this situation and all recharged 
water eventually discharges to the creek. 
The aquifer is unconfined and the groundwater table 
is controlled at the downstream boundary by the water level of the creek.
This is implemented in the model through a specified head boundary. 
The upstream end of the model, i.e., the groundwater divide is implemented
through a no flow boundary. Similarly the base of the aquifer and the
vertical boundary at the creek's position are implemented 
as no flow boundaries. The completed groundwater flow model 
will be provided to you and form the starting point of this exercise.

Save the files supplied for this exercise 
(ORTI file and database file) in a newly created 
folder. Then open the file  {\color{red}\textbf{dd\_u.orti}}.
Before running the model, inspect the model and check
\begin{itemize}
\item the position of the boundary condition under 
{\color{blue}Spatial attributes} 
$\rightarrow$ {\color{blue}\textbf{Modflow}} 
$\rightarrow$ {\color{blue}\textbf{Bas.3}}
\item the initial and specified heads 
{\color{blue}\textbf{Spatial attributes}} 
$\rightarrow$ {\color{blue}\textbf{Modflow}} 
$\rightarrow$  {\color{blue}\textbf{Bas.5}}
\item the implementation of groundwater recharge under 
{\color{blue}\textbf{Spatial attributes}} 
$\rightarrow$ {\color{blue}\textbf{Modflow}} 
$\rightarrow$ {\color{blue}\textbf{Rch.1}}
\item the length of the simulation under 
{\color{blue}\textbf{Parameters}} 
$\rightarrow$ {\color{blue}\textbf{Model}}
$\rightarrow$ {\color{blue}\textbf{Time}}
\end{itemize}

In the next step write the MODFLOW input files and run MODFLOW.
Verify that the {\color{blue}\textbf{Total Time}} shown in the  
pop-up window is the same as your user-specified simulation time. 
If it differs the model run was
not successfully terminated and the model needs to be debugged. 
Inspect the simulation results.

\subsection{Conservative solute transport in a single domain}

Before simulating the reactive transport of uranium we will 
first investigate the transport behavior of a conservative tracer
between a pit that is filled with tailings material and the point of
groundwater discharge to the surface water.
This will allow a first estimate of the travel-time and of the degree
to which mixing and dispersion causes a concentration attenuation 
for any solutes that are transported between the pit and the discharge zone.
In a first step the simulation will be performed under the assumption that 
the aquifer can be described as a media behaving equivalent to a porous media. 
 For these first model simulation the tracer is added across the whole volume 
of water representing the tailings fluids. \\
To run this first conservative transport model follow these steps:

\begin{itemize}
\item set the duration of the simulation (parameters/model/time) to 100 years 
(translate in days)
\item	set the position of the pit for the initial concentration 
of the tracer. We will suppose that the initial concentration 
is homogeneous over the whole pit. 
To do this go to 
{\color{blue}\textbf{Spatial attributes}} 
$\rightarrow$ {\color{blue}\textbf{Mt3dms}} 
$\rightarrow$ {\color{blue}\textbf{Btn.13}}
concentrations and create a zone, that extends from 501 to 799 m in x direction 
and 60 to 89 m in z direction. Set a value of 10 for the tracer. 
This can be any unit.
\item Set the porosity to a low value (0.02) corresponding to the fractures, in 
{\color{blue}\textbf{Spatial attributes}} 
$\rightarrow$ {\color{blue}\textbf{Mt3dms}} 
$\rightarrow$ {\color{blue}\textbf{Btn.11}}
\item Set the longitudinal dispersivity to 1m 
under 
{\color{blue}\textbf{Parameters}} 
$\rightarrow$ {\color{blue}\textbf{Transport}} 
$\rightarrow$ {\color{blue}\textbf{Parameters}}
$\rightarrow$ {\color{blue}\textbf{Dsp.2}}
and leave the ratio of transverse to longitudinal at it default value.
\item Then write the Mt3dms input files and run the model
\end{itemize}

After the model run, look at the temporal evolution of the concentrations 
in the domain to assess the time it takes to deplete the tracer from the pit. 
To view the potential tracer output to the creek, go to 
{\color{blue}\textbf{Results}} 
$\rightarrow$ {\color{blue}\textbf{Observations}}
$\rightarrow$ {\color{blue}\textbf{Breakthrough}}
$\rightarrow$ {\color{blue}\textbf{Creek$\_$obs}}
and select {\color{blue}\textbf{Value}}. 
This will plot the breakthrough curve of the tracer at the point of 
discharge to the creek. Export the results to a text file 
such that the data can later be used for a comparison of different cases.

\subsection{Conservative solute transport in a dual-domain system}

Up to now groundwater flow was considered to occur 
more or less homogeneously in a porous medium. 
However, in fractured rock the vast majority of the water flux occurs
in the fractures while the water prevailing outside the 
fractures can be considered as immobile.
The exchange of solutes between the mobile fraction of the aquifer, 
i.e., the fractures and the immobile domain (i.e., matrix) occurs slowly.
In many cases mass exchange is limited to transport by molecular diffusion. 
This can be modelled by invoking the {\color{blue}\textbf{RCT}} package in MT3DMS. 
In the {\color{blue}\textbf{RCT}} package two major adjustments have to be made,
i.e., the total (single) porosity has to be replaced by
a combination of a mobile and an immobile porosity.
Additionally the exchange coefficient between the mobile domain 
and the immobile domain has to defined. 

\begin{itemize}
\item First add the {\color{blue}\textbf{RCT}} package to the suite of 
modules that can be used in your simulation. To do this go to 
{\color{blue}\textbf{Menu}} 
$\rightarrow$ {\color{blue}\textbf{Addins}},
select {\color{blue}\textbf{Mt3dms modules}} 
and check the {\color{blue}\textbf{RCT}} package. 
This will enable the {\color{blue}\textbf{Spc}} 
and the {\color{blue}\textbf{Rct}} buttons
to appear under 
{\color{blue}\textbf{Parameters}} $\rightarrow$ {\color{blue}\textbf{Transport}}.
\item Go to 
{\color{blue}\textbf{Parameters}}
$\rightarrow$ {\color{blue}\textbf{Transport}} 
$\rightarrow$ {\color{blue}\textbf{Parameters}}
$\rightarrow$ {\color{blue}\textbf{RCT}}
$\rightarrow$ {\color{blue}\textbf{Rct.1 Major Flags}}
and set the {\color{blue}\textbf{ISOTHM}} 
value to {\color{blue}\textbf{dual domain no sorption}}. 
Furthermore, keep 
{\color{blue}\textbf{no reaction}} and set the 
{\color{blue}\textbf{initial conc}} to {\color{blue}\textbf{Given}}, 
before clicking {\color{blue}\textbf{Apply}}.
\item Then, in the same panel, go to 
{\color{blue}\textbf{Rct.2b}} and set the {\color{blue}\textbf{immobile porosity}} 
to 0.23, the value that represents the porosity of the immobile domain. 
Click {\color{blue}\textbf{Apply}}
\item Also in the same panel go to {\color{blue}\textbf{Rct.2c}} 
and set the {\color{blue}\textbf{initial concentrations}} to 0 to define that
the initial concentration of the tracer in the matrix is null.
\item We have to specify which species will be submitted 
to dual domain exchange. 
To do this click on the 
{\color{blue}\textbf{Parameters}}
$\rightarrow$ {\color{blue}\textbf{Transport}} 
$\rightarrow$ {\color{blue}\textbf{Spc}}
button, and in the dialog select {\color{blue}\textbf{Mt3dms}} 
and add {\color{blue}\textbf{Tracer}} as the species name.
\item The exchange coefficient value is provided by using the 
{\color{blue}\textbf{Parameters}}
$\rightarrow$ {\color{blue}\textbf{Transport}} 
$\rightarrow$ {\color{blue}\textbf{Rct}}
button, which opens a grid dialog that only shows {\color{blue}\textbf{Tracer}} 
in this incidence. The mass transfer coefficient which controls the
rate at which mass is exchanged between the mobile and immobile
domain can be entered in the {\color{blue}\textbf{SP2}} column. 
Try a low value such as 
$5.10 \times 10^{-06}$ $d^{-1}$. 
The exchange coefficient employed here is set to a small value
because the matrix porosity, i.e., the water-filled pore space 
not occupied by the fractures is made of very small pores and fissures.
However, over long transport timescales such as in the 
present case the effect of matrix diffusion can still be significant.
\item Finally write the MT3DMS input files and run MT3DMS. 
\end{itemize}

To compare your results to the previous single domain simulation 
display the breakthrough curves again (i.e., {\color{blue}\textbf{Tracer}}
for the same observation point {\color{blue}\textbf{creek$\_$obs}} as above), 
export the concentration and compare the two breakthrough curves in excel. 

\\
If you have time
\begin{itemize}
\item Think about what are the major differences between the simulated 
breakthrough curves ?
\item Rerun the model with with different mass transfer coefficients, 
For example, run a variant with a 10-times higher mass transfer coefficient
and another variant with a 10-times lower mass transfer coefficient. 
\item Interpret the differences between the different types
of breakthrough behavior that is generated by the different model variants.
\end{itemize}

\subsection{Multi-component transport in a dual-domain system}

Following from the (conservative) tracer transport simulations
the next step in this exercise is to expand the model to jointly simulate 
the transport of all major groundwater constituents that may play a role
for the evolution of the hydrochemical patterns in the aquifer.  
Initially we will not consider any water rock interactions and the 
first multi component simulation will only consider aqueous 
complexation reactions. For simplicity we consider only two types 
of water compositions. 

First we consider the ambient groundwater composition, to which we refer 
to as {\color{blue}\textbf{background solution}}. 
Secondly, we consider the composition 
of the tailings fluids, which we refer 
to as (i.e., {\color{blue}\textbf{Solution 1}}).
For simplicity we assume that the tailings fluids have the same composition
throughout the entire pit. Both compositions are listed in Table \ref{ex_dd_u}.
\\
To establish the first multi-component simulation

\begin{itemize}
\item import the database file that is provided for this example by going to
{\color{blue}\textbf{Parameters}}
$\rightarrow$ {\color{blue}\textbf{Chemistry}} 
$\rightarrow$ {\color{blue}\textbf{ImpDb - Import Database}}. 
This database includes most of the important uranium species 
that could be relevant for this type of problem. 
\item go to the
{\color{blue}\textbf{Parameters}}
$\rightarrow$ {\color{blue}\textbf{Chemistry}} 
$\rightarrow$ {\color{blue}\textbf{Chemistry}} 
button to enter the two aqueous solutions compositions 
listed in Table \ref{ex_dd_u}. The corresponding aqueous components 
need to be activated. Note, that all components that have an
initial concentration of zero will also need to be activated (i.e., ticked)
as they can be produced by geochemical reactions during the simulation.
\end{itemize}

The next step will be to assign the two defined aqueous solutions, 
i.e., {\color{blue}\textbf{background solution}} 
and {\color{blue}\textbf{Solution 1}}
as initial and boundary conditions. 
By default the {\color{blue}\textbf{background solution}} is used to
assign the initial concentrations throughout the entire model domain.
However, {\color{blue}\textbf{Solution 1}} has to be assigned 
to the specific zone in teh model that represents 
the location of the pit, which contains 
the tailings fluids. 
We will use the same zone for the pit as the one we used in 
the conservative transport simulations. 
\\
In order to do this, 
\begin{itemize}
\item go to 
{\color{blue}\textbf{Spatial attributes}} 
$\rightarrow$ {\color{blue}\textbf{Mt3dms}}
$\rightarrow$ {\color{blue}\textbf{btn.13}} 
and select in the 
{\color{blue}\textbf{Modify zone}} panel, 
the zone that you created for the pit. 
Click on the {\color{blue}\textbf{Value}} button, 
and copy the coordinates 
(using {\color{blue}\textbf{Ctrl+C}}). 
\item Then go to 
{\color{blue}\textbf{Spatial attributes}} 
$\rightarrow$ {\color{blue}\textbf{Pht3d}}
$\rightarrow$ {\color{blue}\textbf{Ph.3}}
and select {\color{blue}\textbf{Initial chemistry}} 
to create a new rectangular zone that represents
the location of the pit. 
Once the dialog for the zone characteristics opens, 
go to the input zone for the zone coordinates 
and paste the previously saved 
coordinates (using {\color{blue}\textbf{Ctrl +V}}) 
to this location. 
Note, that pasting only works if the target cells are 
selected and not the cell content. Before leaving this Tab 
verify that your coordinates are copied correctly. 
Should cut and paste not work properly you can simply manually
modify the coordinates to those used in the tracer simulation. 
For the {\color{blue}\textbf{zone value}}, 
assign the value {\color{blue}\textbf{1.0.0.0}}, 
which means you assign the water composition
defined defined {\color{blue}\textbf{Solution 1}} to this zone, together with
the Phase/Mineral assemblage {\color{blue}\textbf{0}}, 
Exchanger {\color{blue}\textbf{0}}
and Surface {\color{blue}\textbf{0}}. 
All of the latter contain any activated 
entities and therefore no water-rock or water-sediment interactions 
are activated.   
The defined concentrations will be used as initial concentrations 
at the start of the simulation.
\end{itemize}

Generally the physical transport parameters to be 
used in the multi-component PHT3D 
simulation (such as the longitudinal dispersivity) 
will be automatically adopted from the previously 
defined MT3DMS input. However, some modifications
have to be made with respect to the settings that control
the dual-domain transport behavior.
First, we need to specifically list all aqueous components that
will be affected by the dual-domain mass transfer process
and secondly we need to define the mass transfer coefficient 
for each of the aqueous components. 
With rare exemptions all the aqueous components that are 
included in the reaction network will be affected by the dual-domain
mass transfer.
\\
In order to define the dual-domain mass transfer for the 
multi-component simulation 
\begin{itemize}
\item First click on the 
{\color{blue}\textbf{Parameters}} 
$\rightarrow$ {\color{blue}\textbf{Transport}}
$\rightarrow$ {\color{blue}\textbf{Spc}}
button and specify {\color{blue}\textbf{Pht3d}} instead 
of the previously used {\color{blue}\textbf{Mt3dms}}. 
Then in the dialog box you will need to enter the names 
of all components that are included in the reaction network
Note, that the spelling has to be identical with
the spelling used in the database. You can, for example,
prepare the list in excel or notepad
and then copy and paste it into the dialog box.
\item Then click on the {\color{blue}\textbf{Rct}} 
button that is located next to the
{\color{blue}\textbf{Spc}} button. You shall see there all the 
names of the components that you just entered. 
Then in the {\color{blue}\textbf{SP2}} column specify 
the exchange coefficient to a value of 
$5.10 \times 10^{-06}$ $d^{-1}$
for all listed components, 
including {\color{blue}\textbf{pH}} and {\color{blue}\textbf{pe}}.
\end{itemize}

Once this input is completed the PHT3D input files can be 
written and PHT3D can be executed. 
You will notice that the PHT3D simulation will proceed 
much more slowly than the MT3DMS simulation.
The reason is that PHT3D has to compute the solute transport for 
all included aqueous components and additionally 
the geochemical reactions.
So it is generally wise to carefully verify the model input 
before running the model. If required, it is possible 
to stop the simulation with {\color{blue}\textbf{Ctrl+C}}.
If needed, the intermediate results that were computed up to the last 
simulation time can then be visualised. This can be useful
when a first quick check is required. 

Once the model simulation is terminated 
inspect the transport behaviour of uranium.
To do this visualise the 2D plume contours and 
plot the breakthrough curve.

\begin{itemize}
\item What is the maximum uranium concentration discharging to the creek ?
\item Compare the arrival time of the uranium peak with the peak of 
the previous MT3DMS {\color{blue}\textbf{Tracer}} simulation.        
\end{itemize}


\begin{table}[hb]
\caption{Aqueous concentrations of the ambient groundwater (Background solution)
and of the tailings fluid (Solution 1).}
\begin{center}
\begin{tabular}{l c c}
\hline
 Aqueous component       & $C_{Background}$         & $C_{Solution 1}$      \\
                         & $(mol\ l^{-1}_w)$        & $(mol\ l^{-1}_w)$      \\ \hline
 pH                      &   6.67                   &  4.70            \\
 pe                      &   4                      &  4            \\
 Na                      & 6.16 $\times$ $10^{-5}$  & 5.24 $\times$ $10^{-3}$ \\ 
 Cl                      & 6.13 $\times$ $10^{-3}$  & 1.48 $\times$ $10^{-2}$ \\
 K                       & 5.12 $\times$ $10^{-5}$  & 4.61 $\times$ $10^{-3}$ \\
 Ca                      & 1.15 $\times$ $10^{-3}$  & 5.23 $\times$ $10^{-3}$ \\
 Mg                      & 5.50 $\times$ $10^{-3}$  & 4.44 $\times$ $10^{-1}$ \\
 Al                      & 1.00 $\times$ $10^{-4}$  & 5.23 $\times$ $10^{-3}$ \\
 Amm                     & 0                        & 2.85 $\times$ $10^{-2}$ \\
 C(4)                    & 1.88 $\times$ $10^{-3}$  & 1.40 $\times$ $10^{-5}$ \\
 Mn(2)                   & 0                        & 2.86 $\times$ $10^{-2}$ \\
 Mn(3)                   & 0                        & 0 \\
 S(6)                    & 2.94 $\times$ $10^{-3}$  & 4.92 $\times$ $10^{-1}$ \\
 S(-2)                   & 0                        & 0 \\
 Si                      & 4.00 $\times$ $10^{-3}$  & 3.03 $\times$ $10^{-4}$ \\
 Uranyl                  & 8.00 $\times$ $10^{-7}$  & 3.14 $\times$ $10^{-5}$ \\ \hline
\end{tabular}
\label{ex_dd_u}
\end{center}
\end{table}


\subsection{Development of surface complexation model}

\subsubsection{Experimental data}

Cores were collected from a fractured rock aquifer adjacent to a 
uranium mine. The cores were collected near a pit from the mine that had 
been backfilled with wastewater water and waste-rock from processing 
to extract the uranium. Processing the ore involved extraction 
with sulfuric acid and ammonia. 

The specific surface area of the material determined 
by N$_2$ adsorption was 15 m$^2$/g and the sorbed U(VI) 
concentration on the material determined by carbonate 
extraction \citep{KohlerMatthias2004MfEA} was 80 nmol/g.

The compositions of the groundwater from the fracture zone 
and of water from the pit were determined and used to 
prepare artificial solutions for conducted U(VI) sorption-desorption 
experiments.  The pit-water and groundwater were considered to 
represent end members of the solution compositions that would 
result from pit-water entering the fracture zone.  
Experiments were conducted in artificial pit-water, 
artificial groundwater, and two other groundwater compositions 
prepared based on mixing calculations of the two end members.  
The end-member artificial groundwater was designated AGW1, 
the artificial pit-water APW, the two mixtures AGW3 and AGW4.  
Compositions were presented in Table \ref{ex_u_scm}

\begin{table}[hb]
\caption{Synthetic solutions used in U(VI) sorption-desorption experiments.  All concentrations are in millimoles per liter except pH (NIST pH scale).}
\begin{center}
\begin{tabular}{l c c c c}
\hline
 Aqueous       & $C_{AGW1}$              & $C_{APW}$               & $C_{AGW3}$              & $C_{AGW4}$ \\
 component     & $(mol\ l^{-1}_w)$       & $(mol\ l^{-1}_w)$       &   $(mol\ l^{-1}_w)$     & $(mol\ l^{-1}_w)$ \\ \hline
 pH            &   6.64                  &  4.70                   & 6.50                    & 6.76 \\
 Alkalinity    & 4.05 $\times$ $10^{-3}$ & 0                       & 3.65 $\times$ $10^{-3}$ & 3.93 $\times$ $10^{-3}$ \\ 
 Na            & 1.26 $\times$ $10^{-4}$ & 5.00 $\times$ $10^{-3}$ & 6.13 $\times$ $10^{-4}$ & 2.72 $\times$ $10^{-4}$ \\ 
 K             & 5.10 $\times$ $10^{-5}$ & 4.40 $\times$ $10^{-3}$ & 4.86 $\times$ $10^{-4}$ & 1.82 $\times$ $10^{-4}$ \\ 
 Mg            & 1.97 $\times$ $10^{-3}$ & 3.00 $\times$ $10^{-1}$ & 3.18 $\times$ $10^{-2}$ & 1.09 $\times$ $10^{-2}$ \\
 Ca            & 1.15 $\times$ $10^{-4}$ & 5.00 $\times$ $10^{-3}$ & 6.04 $\times$ $10^{-4}$ & 2.62 $\times$ $10^{-4}$ \\
 Amm           & 0                       & 3.50 $\times$ $10^{-2}$ & 3.50 $\times$ $10^{-3}$ & 1.05 $\times$ $10^{-3}$ \\
 Mn(2)         & 0                       & 2.70 $\times$ $10^{-2}$ & 2.70 $\times$ $10^{-3}$ & 2.00 $\times$ $10^{-5}$ \\
 Zn            & 0                       & 1.50 $\times$ $10^{-4}$ & 1.50 $\times$ $10^{-5}$ & 5.00 $\times$ $10^{-6}$ \\
 Uranyl        & 8.00 $\times$ $10^{-7}$ & 3.08 $\times$ $10^{-5}$ & 3.80 $\times$ $10^{-6}$ & 1.80 $\times$ $10^{-6}$ \\
 Cl            &  0                      & 1.00 $\times$ $10^{-2}$ & 1.00 $\times$ $10^{-3}$ & 3.00 $\times$ $10^{-4}$ \\
 N(5)          &  0                      & 6.00 $\times$ $10^{-5}$ & 6.00 $\times$ $10^{-6}$ & 2.00 $\times$ $10^{-6}$ \\ 
 S(6)          & 1.50 $\times$ $10^{-4}$ & 3.49 $\times$ $10^{-1}$ & 3.50 $\times$ $10^{-2}$ & 9.80 $\times$ $10^{-3}$ \\
 P$_{CO_2}$    & 3.50 $\times$ $10^{-5}$ & 4.00 $\times$ $10^{-4}$ & 3.50 $\times$ $10^{-5}$ & 3.50 $\times$ $10^{-5}$ \\ \hline
\end{tabular}
\label{ex_u_scm}
\end{center}
\end{table}

All experiments were conducted at a solid/liquid ratio of 10 g/L at 
25 degrees C. Experiments with AGW1, AGW3, and AGW4 were conducted 
so that equilibrium with a P$_{CO_2}$ of 0.035 atmospheres was maintained.  
Experiments with APW were conducted in equilibrium with the 
atmosphere (0.00040 atmosphere).  
After steady concentrations of U(VI) were observed, 
the suspensions were considered to be at equilibrium.

The total U(VI) concentration in each experiment was calculated from the U(VI) 
added in the aqueous phase and the sorbed U(VI) added with the solid phase.  
Total U(VI) concentrations and the sorbed U(VI) concentrations 
at equilibrium are presented in Table \ref{ex_u_scm2}.

\begin{table}[h]
\caption{Total U(VI) concentrations and sorbed U(VI) concentrations 
at equilibrium in each of the solution used in the experiments.}
\begin{center}
\begin{tabular}{l c c c}
\hline
 Solution                & Total U(VI)              & Dissolved U(VI)         & Sorbed U(VI)        \\
                         & $(mol\ l^{-1}_w)$        & $(mol\ l^{-1}_w)$       & $(mol\ l^{-1}_w)$   \\ \hline
 AGW1                    & 8.00 $\times$ $10^{-7}$  & 1.90 $\times$ $10^{-7}$ & 6.10 $\times$ $10^{-7}$ & \\
 APW                     & 3.08 $\times$ $10^{-5}$  & 2.16 $\times$ $10^{-5}$ & 1.06 $\times$ $10^{-5}$ & \\ 
 AGW3                    & 3.80 $\times$ $10^{-6}$  & 1.40 $\times$ $10^{-6}$ & 2.40 $\times$ $10^{-6}$ & \\
 AGW4                    & 1.80 $\times$ $10^{-6}$  & 1.25 $\times$ $10^{-6}$ & 5.50 $\times$ $10^{-7}$ & \\ \hline
\end{tabular}
\label{ex_u_scm2}
\end{center}
\end{table}

\begin{table}[h]
\caption{Possible surface complexation reactions to fit the experimental data.}
\begin{center}
\begin{tabular}{l c c c}
\hline
 Surface reaction           \\ \hline
 $SOH$ + $UO_2^{2+}$ + $H_2O$ + $CO_2$ = $SOHUO_2HCO_3+$ + $H^+$  \\
 $SOH$ + $UO_2^{2+}$ + $H_2O$ + $CO_2$ = $SOUO_2H_2CO_3$ + $2H^+$  \\ \hline
\end{tabular}
\label{ex_u_scm3}
\end{center}
\end{table}

\subsubsection{Site-specific surface complexation model}

To describe the experimental data perform the following steps
\begin{itemize}
\item Choose one of the two surface complexation reactions 
in Table \ref{ex_u_scm3}. 
\item Fix the total site concentration at 5.76 $\times$ $10^{-4}$.  
\item Use a non-electrostatic model. 
\item Fit the sorbed U(VI) concentration in Table \ref{ex_u_scm2} 
by adjusting the logK value.   
\item Find the logK value that you think provides the closest 
fit to the experimental data. 
\item You may use the provided Phreeqc input and the provided Excel file
\end{itemize}

Note, that neither of the reactions in Table \ref{ex_u_scm3} 
will provide a perfect fit to the experimental data

\subsection{Impact of surface complexation reactions 
in the field-scale model}

In the next step we will investigate how sorption of uranium affects
the overall uranium migration patterns. To do this we will activate the
surface complexation model that was developed for this site. 

\begin{itemize}
\item Go to 
{\color{blue}\textbf{Parameters}}
$\rightarrow$ {\color{blue}\textbf{Chemistry}} 
$\rightarrow$ {\color{blue}\textbf{Chemistry}}
and select {\color{blue}\textbf{Surface}}.
Now, activate the surface with the name {\color{blue}\textbf{Surf\_a}} 
and enter 2.0 $\times$ $10^{-4}$ as the surface site
concentration under {\color{blue}\textbf{Site\_back}}.
This will cause the aqueous uranium that migrates away
from the pit to undergo surface complexation reactions,
with the variable extent of the sorption being a function 
of the spatially and temporally variable geochemical conditions.
\item The present model is using the {\color{blue}\textbf{-no\_edl}} 
option, which means that the surface complexation reactions
do not consider any electric double layer effects.
If required, this can be changed under 
{\color{blue}\textbf{Parameters}}
$\rightarrow$ {\color{blue}\textbf{Chemistry}} 
$\rightarrow$ {\color{blue}\textbf{Parameters}}
$\rightarrow$ {\color{blue}\textbf{PH}}
$\rightarrow$ {\color{blue}\textbf{ph.6}}
where a different option such  {\color{blue}\textbf{-ddl}} 
could be entered. 
\end{itemize}

Now rerun PHT3D and investigate how much the sorption reactions have impacted 
the uranium transport behaviour.   

\subsection{Impact of mineral buffering}

So far we have excluded mineral reactions and therefore not considered 
the potential impact of pH buffering reactions on 
raising the pH of the aqueous solution that emerges from the 
tailings pit zone. In the next step we add chlorite to the 
background mineral assemblage. 
\\
To implement chlorite
\begin{itemize}
\item Go to 
{\color{blue}\textbf{Parameters}}
$\rightarrow$ {\color{blue}\textbf{Chemistry}} 
$\rightarrow$ {\color{blue}\textbf{Chemistry}}
and select {\color{blue}\textbf{Phase}}.
Now, activate chlorite and enter a background concentration
of 0.1 mol/L under {\color{blue}\textbf{Backgr. Moles}}, i.e., 
in the second column.
\item Then go to 
{\color{blue}\textbf{Spatial attributes}} 
$\rightarrow$ {\color{blue}\textbf{Pht3d}}
$\rightarrow$ {\color{blue}\textbf{Ph.3}}
and select {\color{blue}\textbf{Initial chemistry}} 
to define that chlorite is not present in the pit zone.
Change the previously entered value 
{\color{blue}\textbf{1.0.0.0}} to {\color{blue}\textbf{1.1.0.0}}.
This means that you now assign 
the Phase/Mineral assemblage {\color{blue}\textbf{1}}
to this zone. As the chlorite concentration for Mineral assemblage 1 is zero
the initial chlorite concentration that will be used for the pit zone
will also be zero.
\end{itemize}.

Now rerun the model and investigate the impact on the uranium transport behavior.


 
  




